\begin{figure}[t]
\centering
\begin{tikzpicture}[
 box/.style={draw=blue!35!black,rounded corners=2pt,fill=blue!3,
  text width=7.2cm,align=center,inner sep=5pt,font=\small},
 branch/.style={box,text width=5cm},
 flow/.style={-{Latex[length=2mm]},semithick,draw=blue!50!black}]
\node[box] (spec) {GOE edge and resolvent control\\Proposition~\ref{prop:spectral}};
\node[box,below=8mm of spec] (sphere) {Spherical estimates from Gaussian conditioning\\Partition, projection, and overlap bounds\\Section~\ref{sec:spherical}};
\node[box,below=8mm of sphere] (cube) {Haar first and second moments\\Partition concentration and observable transfer\\Section~\ref{sec:orientation}};
\node[branch,anchor=north] (small) at ([xshift=-2.8cm,yshift=-12mm]cube.south) {Small-projection mass on the cube\\Theorem~\ref{thm:static-small-ball}};
\node[branch,anchor=north] (cov) at ([xshift=2.8cm,yshift=-12mm]cube.south) {Squared-overlap control on the cube\\Proposition~\ref{prop:overlap}};
\node[branch,below=8mm of small] (slow) {Sign test for realized initial states\\Theorem~\ref{thm:gibbs-start}};
\node[branch,below=8mm of cov] (lin) {Covariance scale and linear\\inverse Rayleigh bounds\\Theorem~\ref{thm:linear-scale}};
\draw[flow] (spec)--(sphere);
\draw[flow] (sphere)--(cube);
\draw[flow] (spec.west)--++(-6mm,0)|-(cube.west);
\draw[flow] (cube.south)--++(0,-.6)-|(small.north);
\draw[flow] (cube.south)--++(0,-.6)-|(cov.north);
\draw[flow] (small)--(slow);
\draw[flow] (cov)--(lin);
\draw[flow] (small.east)--(lin.west);
\end{tikzpicture}
\caption{The lower and static argument. The direct arrow from the spectral estimates to Haar comparison records their additional use in the cubic overlap deficit. The restricted first moment transfers the small-projection numerator, while partition concentration controls its normalization; a weighted second moment transfers squared overlap. Small projections give the covariance lower bound and squared overlap gives its upper bound. A separate cavity estimate then controls the linear inverse Rayleigh quotients. The independent subexponential upper comparison is outside the diagram.}
\label{fig:proof-structure}
\end{figure}
